\subsection{Элементарная структура базового дифференцирования}
Работаем над $K=\mathbb Q[a]/(a^2-a+1)$ и полагаем
\[
b=1-a,\qquad v=z+axy,\qquad
A(x)=-1+ax^2-ax^4,
\]
\[
q(x)=-2x^5+(6-2a)x^3-2bx,
\qquad h_0(x,v)=bxv^2+A(x)v+q(x).
\]
Прямое раскрытие скобок даёт более короткую запись базовой координаты:
\[
H(x,y,z)=h_0(x,z+axy)-bx^5y.
\]
Замена $(x,y,z)\leftrightarrow(x,v,y)$ полиномиальна в обе стороны.
Якобиево дифференцирование
\[
D=H_z\frac{\partial}{\partial y}-H_y\frac{\partial}{\partial z}
\]
в этих координатах имеет треугольный вид
\[
\boxed{D=bx^5\frac{\partial}{\partial v}
       +(2bxv+A(x))\frac{\partial}{\partial y}.}
\]
Здесь $D(x)=D(H)=0$, $D^2(v)=0$, $D^3(y)=D^3(z)=0$;
следовательно, $D$ локально нильпотентно. Это утверждение следует
не из численной проверки, а из указанного треугольного вида.

Пусть $f(x,T)$ --- полином из репозитория. Точные вычисления дают
\[
f(x,h_0(x,v))\equiv v\pmod{x^5},\qquad
h_0(x,f(x,T))\equiv T\pmod{x^5}.
\]
Эти тождества можно также получить конечным обращением
$h_0(x,v)\equiv-v\pmod{x}$ над $K[x]/(x^5)$.
Поэтому выражение
\[
W=\frac{a\bigl(v-f(x,H)\bigr)}{x^5}
\]
полиномиально. Так как $D(v)=bx^5$, $D(f(x,H))=0$ и $ab=1$, имеем
\[
\boxed{D(W)=1.}
\]
В частности, $W$ --- явный слайс. Для полноты координатности
предъявляется обратное преобразование:
\[
v=f(x,T)+bx^5V,\qquad
y=\frac{h_0(x,v)-T}{bx^5},\qquad z=v-axy.
\]
Последнее деление точно вследствие второго сравнения выше.
Эти формулы взаимно обратны к $(x,y,z)\mapsto(x,H,W)$ после
обращения $x$. Поскольку обе стороны полиномиальны, инъективность
локализации устанавливает обе композиции и на $x=0$.
Таким образом,
\[
K[x,y,z]=K[x,H,W],\qquad \ker D=K[x,H].
\]
Это конкретная реализация классического обращения по модулю
исключительного множителя; общая лемма не заявляется как новая.

Для гармонического переноса положим
\[
h_H=H(x,xz-y^2+2,y),\qquad
D_H=(h_H)_z\partial_y-(h_H)_y\partial_z.
\]
Если $q_H$ --- третий компонент относительного автоморфизма
репозитория, то его якобиан равен $+1$, поэтому
$D_H(q_H)=-1$. Следовательно, $-q_H$ есть слайс и
$\ker D_H=K[x,h_H]$. Независимое точное дифференцирование даёт
$D_H^3(y)=0$, $D_H^5(z)=0$, причём предыдущие итерации
$D_H^2(y)$ и $D_H^4(z)$ ненулевые. Это два конкретных
полиномиальных локально нильпотентных дифференцирования с
полностью явными координатами ядра и слайсами.
